\documentclass[11pt,a4paper]{article}
\usepackage[T1]{fontenc}
\usepackage{lmodern}
\usepackage[margin=27mm]{geometry}
\usepackage{amsmath,amssymb,amsthm,mathtools}
\usepackage{booktabs,array,longtable}
\usepackage{microtype}
\usepackage{xcolor}
\usepackage{enumitem}
\usepackage{hyperref}
\hypersetup{colorlinks=true,linkcolor=blue!45!black,citecolor=blue!45!black,urlcolor=blue!45!black,
 pdftitle={Exact restriction laws and Fourier cancellation on hyperbolic quadrics},
 pdfauthor={ChatGPT; research manuscript prepared for the ProveIt project}}
\setlength{\emergencystretch}{2em}
\setlength{\parskip}{3pt}
\setlist{itemsep=3pt,topsep=5pt}
\newtheorem{theorem}{Theorem}[section]
\newtheorem{proposition}[theorem]{Proposition}
\newtheorem{lemma}[theorem]{Lemma}
\newtheorem{corollary}[theorem]{Corollary}
\theoremstyle{definition}
\newtheorem{definition}[theorem]{Definition}
\newtheorem{question}[theorem]{Question}
\theoremstyle{remark}
\newtheorem{remark}[theorem]{Remark}
\newcommand{\F}{\mathbb F}
\newcommand{\C}{\mathbb C}
\newcommand{\E}{\mathbb E}
\newcommand{\Pp}{\mathbb P}
\newcommand{\Sym}{\operatorname{Sym}}
\newcommand{\rank}{\operatorname{rank}}
\newcommand{\tr}{\operatorname{tr}}
\newcommand{\im}{\operatorname{im}}
\newcommand{\Var}{\operatorname{Var}}
\newcommand{\Cov}{\operatorname{Cov}}
\newcommand{\diag}{\operatorname{diag}}
\newcommand{\GL}{\operatorname{GL}}
\newcommand{\Gr}{\operatorname{Gr}}
\newcommand{\eps}{\varepsilon}
\newcommand{\J}{\mathsf J}
\newcommand{\D}{\mathsf D}
\newcommand{\I}{\mathsf I}
\newcommand{\G}{\mathcal G}
\newcommand{\one}{\mathbf 1}
\newcommand{\gb}[2]{\genfrac{[}{]}{0pt}{}{#1}{#2}_{q}}
\newcommand{\abs}[1]{\left|#1\right|}
\newcommand{\norm}[1]{\left\|#1\right\|}
\title{Exact restriction laws and Fourier cancellation\\on hyperbolic quadrics\\[5pt]
\large Sharp polar averaging and two-family discriminant discrepancies}
\author{Research manuscript prepared with ChatGPT\\for the ProveIt project}
\date{7 October 2026}
\begin{document}
\maketitle
\begin{abstract}
Let $H$ be a split nondegenerate symmetric form on $\F_q^{2r}$, where $q$ is odd,
and let $Q$ be a uniformly random symmetric form. We investigate the determinant
characters of the restrictions of $Q$ to the maximal $H$-isotropic subspaces.
An exact conditional extension law gives the joint sign distribution for two
arbitrary subspaces, including every possible singular common restriction.
For the average over all generators we prove the sharp fixed-rank variance
asymptotic $q^{-(r-1)}(1+O_r(q^{-1}))$ for determinant character.
A different phenomenon occurs on subtracting the averages over the two classical
Latin--Greek families: every Fourier frequency of rank less than $r$ cancels.
The resulting exact second moment is
\[
 \E\Delta_{\D}^{\,2}
 =\frac{4p_r}{N_r}q^{-\lfloor r/2\rfloor},\qquad
 p_r=\prod_{\substack{1\le j\le r\\j\text{ odd}}}(1-q^{-j}),\quad
 N_r=2\prod_{j=1}^{r-1}(q^j+1).
\]
Its exponent is $\lfloor r^2/2\rfloor$, rather than $r-1$.
In even rank the character contrast is, pointwise, a signed count of common
generators times an explicit scalar. It is exactly zero with probability
$1-(2+O_r(q^{-1}))q^{-r}$, and a nonzero value forces a square characteristic
polynomial for $H^{-1}Q$. The variance-normalized contrast converges to zero
in probability, despite having second moment one. An analogous odd-rank
identity holds for nonsingularity.
We also obtain exact nonsingularity and individual-sign contrast variances,
a general incidence-to-Fourier transfer principle, and sharp Grassmannian
averaging. These results yield an $O_r(q^{-1})$ local transverse approximation
for $r\ge3$ in the rank-layer setup of the cited OpenAI manuscript, replacing
its $O_r(q^{-1/2})$ estimate. All theorems have proofs here; exhaustive finite
checks and integration notes accompany the article. Classical rank/type counts,
Fourier ingredients, and the Latin--Greek halving are explicitly distinguished
from the proposed research contribution. Literature priority, independent
referee approval, and proof-assistant verification are not asserted.
\end{abstract}
\newpage
\tableofcontents
\newpage

\section{Research setting, source review, and principal conclusions}
\label{sec:intro}

\subsection{A specific local problem selected from the repositories}
The direct mathematical starting point is the section \emph{Transverse
configurations and determinant signs} in the manuscript filed as
\emph{A counterexample to Sidorenko's conjecture} in \texttt{openai/math}
\cite{openai}. Its lemma labelled \texttt{lem:transverse-restrictions} treats
any family $\mathcal U$ of distinct $r$-subspaces and proves
\begin{equation}\label{eq:source-bound}
 \E\abs{\frac1{|\mathcal U|}\sum_{U\in\mathcal U}
 \one_{\{\chi(\det(Q|_U))=\xi\}}-\frac12}
 =O_r\bigl(q^{-1/2}+|\mathcal U|^{-1/2}\bigr),\qquad \xi=\pm1.
\end{equation}
The source uses this for maximal isotropic subspaces of a split form and for
all middle-dimensional subspaces. Its proof bounds every off-diagonal
covariance by $O_r(q^{-1})$, discarding the dimension and nullity of the common
restriction. We retain that information exactly.

The selected ProveIt context is its \emph{Local Quantitative Refinements of
Gowers's Proof of Szemer\'edi's Theorem} program \cite{proveit}. The reviewed
inventory includes sharp moment and density-transfer inequalities, refinements
of restriction estimates, and collision-sensitive improvements of local
counting arguments. Its source-33 audit explicitly separates written proofs,
finite tests, source interfaces, and unresolved priority. This article follows
that local-theorem approach, but does not use an unverified ProveIt theorem as
a premise. The connection to ProveIt is its quantitative counting and
restriction program; the direct external theorem interface comes from
\texttt{openai/math}.

This was a scoped review of repository inventories and selected mathematical
sources, not a proof audit of every manuscript in either repository. In
particular, the title of the Sidorenko manuscript is not treated as a
certification of its global conclusion. Our finite-field theorems stand
independently of that conclusion.

\subsection{Three distinct scales}
Fix an odd prime power $q$ and a split form $H$ on a $2r$-dimensional vector
space. Write $\G$ for its maximal totally isotropic $r$-subspaces, and
$\G_+,\G_-$ for its two equal-sized classical families. For a symmetric form
$C$ set
\[
 \J(C)=\one_{\{\det C\ne0\}},\qquad
 \D(C)=\chi(\det C),\qquad
 \I_\xi(C)=\frac{\J(C)+\xi\D(C)}2.
\]
Thus $\I_\xi$ is the indicator of a nonsingular restriction of sign $\xi$;
$\D$ is zero, not a quotient discriminant, on a singular matrix.
For $f\in\{\J,\D,\I_\xi\}$ define
\begin{equation}\label{eq:averages-intro}
 \overline f(Q)=\frac1{N_r}\sum_{U\in\G}f(Q|_U),\qquad
 F_{f,\pm}(Q)=\frac2{N_r}\sum_{U\in\G_\pm}f(Q|_U),\qquad
 \Delta_f=F_{f,+}-F_{f,-}.
\end{equation}
All expectations below are with respect to uniform $Q$ on the full vector
space of symmetric forms, unless conditioning is explicitly displayed.

Theorems~\ref{thm:polar} and \ref{thm:contrast} prove, for fixed $r\ge2$,
\begin{align}
 \Var(\overline\D)&=(1+O_r(q^{-1}))q^{-(r-1)},\label{eq:intro-average}\\
 \Var(\overline\I_\xi)&=(\tfrac14+o_r(1))q^{-(r-1)},\label{eq:intro-indicator}\\
 \E\Delta_\D^2&=(2+O_r(q^{-1}))q^{-\lfloor r^2/2\rfloor}.
 \label{eq:intro-contrast}
\end{align}
The third line concerns agreement between the two families, not proximity of
either family to an unbiased coin. Their means have a deterministic
finite-field bias, and their much larger common fluctuations mostly cancel
in the difference. In particular, replacing the true mean by $1/2$ before
estimating variance would lose the principal distinction.

\begin{center}
\begin{tabular}{@{}lrrrrrrr@{}}
\toprule
Generator dimension $r$ &2&3&4&5&6&7&8\\
\midrule
All-generator variance exponent &1&2&3&4&5&6&7\\
Two-family character exponent &2&4&8&12&18&24&32\\
\bottomrule
\end{tabular}
\end{center}
These are exponents of $q^{-1}$ in variance, not in standard deviation.

\subsection{Attribution and the claim boundary}
The classification of symmetric forms by rank and determinant type, exact
rank/type counts, and their translation association schemes are classical;
Schmidt gives explicit counts and eigenvalues \cite{schmidt2015,schmidt2020}.
The Fourier calculations below are elementary specializations of that
established theory, not new general eigenvalue formulas. The two generator
families and their equality of proper-subspace incidences are the classical
Latin--Greek halving; see Kiermaier, Schmidt and Wassermann
\cite{polar-designs}, especially their discussion preceding Lemma~14.

The proposed contribution is the assembled restriction-covariance calculus,
its sharp averaging laws, the explicit discriminant contrast identities,
and the resulting local rank-layer improvement. We give direct proofs so
that these statements can be checked without importing the full association
scheme machinery. A literature search sufficient to identify the classical
background is not sufficient to certify that no equivalent specialization
has appeared elsewhere. Accordingly, these are proved research results
presented for review, not a claim to have settled a named longstanding
conjecture or established historical priority.

\section{Symmetric forms: normalization and exact marginal laws}
\label{sec:counts}
Throughout, $q\ge3$ is an odd prime power, $\chi:\F_q\to\{0,1,-1\}$ is the
quadratic character with $\chi(0)=0$, and $\eps=\chi(-1)$. Write
\[
 c(n)=\frac{n(n+1)}2,\qquad \Sym_n=\{C\in M_n(\F_q):C=C^{\mathsf T}\}.
\]
A quadratic form is represented by $x^{\mathsf T}Cx$. Odd characteristic
identifies this description with symmetric bilinear forms; all determinant
characters in this article use the displayed matrix convention. Empty
matrices have determinant $1$. In particular $\J_0=\D_0=1$.

For uniform $C\in\Sym_n$, let
\begin{equation}\label{eq:marginal-defs}
 p_n=\E\J(C),\qquad \mu_n=\E\D(C),\qquad
 \pi_{n,\xi}=\Pp(\D(C)=\xi)=\frac{p_n+\xi\mu_n}{2}.
\end{equation}
For $n=0$ the last formula means $\pi_{0,+}=1$ and $\pi_{0,-}=0$.
The Gaussian coefficient $\gb nk$ counts the $k$-subspaces of $\F_q^n$;
it is zero outside $0\le k\le n$. We use
\begin{equation}\label{eq:G-rho}
 G_{t,k}=\gb tk q^{c(t-k)-c(t)},\qquad
 \rho_{k,a}=\prod_{j=a-k+1}^{a}(1-q^{-j})\quad(0\le k\le a),
\end{equation}
and set $\rho_{k,a}=0$ for $k>a$. Thus $\rho_{k,a}$ is the probability that
a uniform $k$-by-$a$ matrix has full row rank.

\begin{lemma}[Radical and quotient counting]\label{lem:radical}
For uniform $C\in\Sym_t$, the probability that $C$ has nullity $k$ and its
nondegenerate quotient has determinant character $\sigma$ is
\begin{equation}\label{eq:radical-type}
 \Pp(k(C)=k,\sigma(C)=\sigma)
 =G_{t,k}\frac{p_{t-k}+\sigma\mu_{t-k}}2,\qquad \sigma=\pm1.
\end{equation}
The probability of nullity $k$ alone is $G_{t,k}p_{t-k}$.
\end{lemma}
\begin{proof}
Choose the radical $K$, in $\gb tk$ ways. Forms with radical exactly $K$
are precisely the pullbacks of nonsingular forms on $\F_q^t/K$.
In any chosen quotient basis their signed count is
$q^{c(t-k)}\pi_{t-k,\sigma}$. Division by $q^{c(t)}$ gives the result.
Changes of quotient basis multiply determinants by squares and do not
alter the sign. The zero-dimensional quotient is covered by the empty
matrix convention.
\end{proof}

\begin{proposition}[Exact marginal probabilities; classical]\label{prop:marginals}
For every $n\ge0$,
\begin{equation}\label{eq:p-product}
 p_n=\prod_{\substack{1\le j\le n\\j\text{ odd}}}(1-q^{-j}),
\end{equation}
and
\begin{equation}\label{eq:mu-product}
 \mu_{2m}=\eps^m q^{-m}p_{2m},\qquad \mu_{2m+1}=0.
\end{equation}
Moreover $1-p_n\le 9/(8q)$ for $n\ge1$, and
$|\pi_{n,\xi}-1/2|\le q^{-1}$ for $n\ge2$.
\end{proposition}
\begin{proof}
Consider a symmetric extension of an $(n-1)$-dimensional leading block by
one row and column. If the leading block is nonsingular, Schur complementation
makes the last scalar uniform, so its conditional nonsingularity probability
is $1-q^{-1}$ and its conditional determinant-character mean is zero.
If the leading block has nullity one, eliminate its nondegenerate part.
The remaining block has the form
$\left(\begin{smallmatrix}0&b\\b&d\end{smallmatrix}\right)$ and is nonsingular
exactly when $b\ne0$. Its determinant has character $\eps$ in that case.
Nullity at least two cannot be repaired by one new dimension. Lemma~\ref{lem:radical}
therefore gives, for $n\ge2$,
\begin{align}
 p_n&=(1-q^{-1})p_{n-1}
       +q^{-1}(1-q^{-(n-1)})p_{n-2},\label{eq:p-recursion}\\
 \mu_n&=\eps q^{-1}(1-q^{-(n-1)})\mu_{n-2}.
 \label{eq:mu-recursion}
\end{align}
Here $G_{n-1,1}(1-q^{-1})=q^{-1}(1-q^{-(n-1)})$.
The initial values are $p_0=1$, $p_1=1-q^{-1}$, $\mu_0=1$, $\mu_1=0$.
Induction proves the two products: $p_{2m}=p_{2m-1}$ and
$p_{2m+1}=(1-q^{-(2m+1)})p_{2m}$.

The product inequality $1-\prod(1-x_i)\le\sum x_i$ gives
$1-p_n\le q^{-1}/(1-q^{-2})\le9/(8q)$.
For odd $n\ge3$, $\mu_n=0$. For even $n\ge4$,
$|\mu_n|\le q^{-2}$, so
$|\pi_{n,\xi}-1/2|\le9/(16q)+1/(2q^2)<q^{-1}$.
For $n=2$, use $p_2=1-q^{-1}$ and
$|\mu_2|=q^{-1}(1-q^{-1})$ directly.
\end{proof}

\section{An exact extension law and the two-subspace distribution}
\label{sec:overlap}
\begin{lemma}[Extension through an arbitrary singular block]\label{lem:extension}
Fix $C\in\Sym_t$ of nullity $k$ and quotient determinant character $\sigma$.
Let
\[
 M=\begin{pmatrix}C&B\\B^{\mathsf T}&E\end{pmatrix},
\]
where $B\in M_{t,a}(\F_q)$ and $E\in\Sym_a$ are independent and uniform.
Define, for $k\le a$,
\begin{equation}\label{eq:AB-def}
 A_{a,k}=\rho_{k,a}p_{a-k},\qquad
 B_{a,k}=\eps^k\rho_{k,a}\mu_{a-k},
\end{equation}
and define both quantities to be zero when $k>a$. Then
\begin{equation}\label{eq:extension-moments}
 \E[\J(M)\mid C]=A_{a,k},\qquad
 \E[\D(M)\mid C]=\sigma B_{a,k}.
\end{equation}
\end{lemma}
\begin{proof}
Choose coordinates in which $C=C_0\oplus0_k$, with $C_0$ nonsingular.
Eliminating $C_0$ by congruence leaves, apart from this fixed summand,
\begin{equation}\label{eq:extension-block}
 \begin{pmatrix}0_k&Z\\Z^{\mathsf T}&F\end{pmatrix},
\end{equation}
where $Z$ is uniform $k$-by-$a$ and $F$ is uniform symmetric, independently.
Indeed, conditional on the mixed block adjacent to $C_0$, the Schur
complement is a translation of the uniform matrix $E$.
If $Z$ does not have full row rank, a nonzero vector in its left kernel
is radical for \eqref{eq:extension-block}; if $k>a$ this always occurs.
If $Z$ has full row rank, independent invertible changes of its row and
column coordinates make it $(I_k\ 0)$. Congruence then gives $k$ hyperbolic
planes and a residual form $R\in\Sym_{a-k}$. The changes in the bottom
symmetric block are invertible linear transformations and translations,
so $R$ is uniform, also conditional on $Z$ and the eliminated entries.
The determinant character of a nonsingular extension is
$\sigma\eps^k\chi(\det R)$; all determinant factors from changes of basis
are squares. Multiplying the residual expectations by
$\Pp(\rank Z=k)=\rho_{k,a}$ proves the claim.
\end{proof}

\begin{proposition}[Full rank-and-type transition]\label{prop:full-transition}
In Lemma~\ref{lem:extension}, the rank $\ell$ of the residual rectangular
matrix $Z$ has distribution
\begin{equation}\label{eq:rectangular-rank}
 R_{k,a}(\ell)=q^{-ka}\gb k\ell\gb a\ell\,|\GL_\ell(\F_q)|.
\end{equation}
Conditional on $\ell$, the extension is congruent to
\[
 C_0\ \oplus\ H_\ell\ \oplus\ R\ \oplus\ 0_{k-\ell},
 \qquad R\text{ uniform in }\Sym_{a-\ell},
\]
where $H_\ell$ is a sum of $\ell$ hyperbolic planes. Thus, if $R$ has rank
$j$ and quotient sign $\zeta$, the full extension has rank
$t-k+2\ell+j$ and quotient sign $\sigma\eps^\ell\zeta$.
The probability of the residual pair $(j,\zeta)$ is
$\gb{a-\ell}{j}q^{c(j)-c(a-\ell)}\pi_{j,\zeta}$.
\end{proposition}
\begin{proof}
A rank-$\ell$ linear map is specified by its kernel, image, and an
isomorphism between the resulting $\ell$-dimensional spaces. This gives
\eqref{eq:rectangular-rank}. Reduce $Z$ to a block with $I_\ell$ and zeros
and apply the same congruence elimination as in Lemma~\ref{lem:extension}.
The unmatched $k-\ell$ zero rows remain radical, while the unpaired new
coordinates carry a uniform symmetric form. Lemma~\ref{lem:radical} supplies
its rank/type probability. If several $(\ell,j,\zeta)$ yield the same final
rank and sign, their probabilities are added; the statement is an exact
mixture, not a disjoint parametrization by final rank alone.
\end{proof}

\begin{theorem}[Exact joint restriction-sign law]\label{thm:joint}
Let $U,V$ be subspaces of dimensions $r,s$ in any finite-dimensional
$\F_q$-space, with $t=\dim(U\cap V)$, $a=r-t$ and $b=s-t$.
For uniform ambient symmetric $Q$ and $\xi,\eta\in\{\pm1\}$,
\begin{align}
 &\E\bigl[\I_\xi(Q|_U)\I_\eta(Q|_V)\bigr]\notag\\
 &\quad=\frac14\sum_{k=0}^{t}G_{t,k}\Bigl[
 p_{t-k}\bigl(A_{a,k}A_{b,k}+\xi\eta B_{a,k}B_{b,k}\bigr)
 \notag\\[-2pt]
 &\hspace{47mm}
 +\mu_{t-k}\bigl(\xi B_{a,k}A_{b,k}+\eta A_{a,k}B_{b,k}\bigr)
 \Bigr].\label{eq:joint-law}
\end{align}
The conventions above cover $t=0$, coincident subspaces, and
zero-dimensional spaces without separate limiting arguments.
\end{theorem}
\begin{proof}
Choose an ambient basis beginning with $U\cap V$ and then the separate
extensions to $U$ and $V$. The two restrictions share precisely the leading
symmetric block $C=Q|_{U\cap V}$. Every other entry on which either
restriction depends is independent and uniform, so the restrictions are
conditionally independent given $C$.
For nullity $k$ and quotient sign $\sigma$ of $C$, Lemma~\ref{lem:extension}
gives conditional sign probabilities
\[
 \frac{A_{a,k}+\xi\sigma B_{a,k}}2,
 \qquad \frac{A_{b,k}+\eta\sigma B_{b,k}}2.
\]
Multiply them, weight by \eqref{eq:radical-type}, and sum over
$\sigma=\pm1$. The even terms sum to $p_{t-k}$ and the odd terms to
$\mu_{t-k}$, giving \eqref{eq:joint-law}.
\end{proof}

\begin{corollary}[Exact covariance kernels]\label{cor:kernels}
For equal dimensions $r=s$, put $a=r-t$ and define
\begin{align}
 K_\J(r,t)&=\sum_{k=0}^{t}G_{t,k}p_{t-k}A_{a,k}^{\,2}-p_r^2,
 \label{eq:KJ}\\
 K_\D(r,t)&=\sum_{k=0}^{t}G_{t,k}p_{t-k}B_{a,k}^{\,2}-\mu_r^2,
 \label{eq:KD}\\
 K_{\J\D}(r,t)&=\sum_{k=0}^{t}G_{t,k}\mu_{t-k}A_{a,k}B_{a,k}-p_r\mu_r.
 \label{eq:KJD}
\end{align}
These are respectively the covariances of the two nonsingularity indicators,
the two determinant characters, and the nonsingularity indicator on $U$
with the determinant character on $V$. In particular
\begin{equation}\label{eq:KI}
 K_{\I_\xi}(r,t)=\tfrac14\bigl(K_\J(r,t)+K_\D(r,t)
                          +2\xi K_{\J\D}(r,t)\bigr).
\end{equation}
All three kernels vanish at $t=0$. Also $K_\J,K_\D\ge0$, and
$|K_{\J\D}|\le\sqrt{K_\J K_\D}$.
\end{corollary}
\begin{proof}
Condition on the common restriction as above. For equal dimensions the
conditional expectations are the same functions of $C$, namely
$A_{a,k}$ and $\sigma B_{a,k}$. Their variances and covariance are precisely
the displayed kernels. This proves positivity and the Cauchy--Schwarz
bound, including the case $U=V$. For $t=0$ the common block is deterministic.
\end{proof}

\section{Uniform covariance estimates and intersection profiles}
\label{sec:bounds}
The exact kernels already give an algorithm for every finite parameter.
The following bounds identify which intersection dimensions dominate.
The constants are deliberately simple and uniform in the dimension.

\begin{lemma}[Uniform kernel bounds]\label{lem:kernel-bounds}
For $r\ge1$, $0\le t\le r$ and $a=r-t$,
\begin{equation}\label{eq:kernel-bounds}
 0\le K_\D(r,t)\le4q^{-a},\qquad
 0\le K_\J(r,t)\le64q^{-(2a+1)}.
\end{equation}
Both kernels are exactly zero when $t=0$. For fixed $r,t$ with $t\ge1$,
\begin{equation}\label{eq:KD-asymptotic}
 K_\D(r,t)=q^{-(r-t)}(1+O_r(q^{-1})).
\end{equation}
\end{lemma}
\begin{proof}
First note the dimension-free estimates
\begin{equation}\label{eq:gaussian-bounds}
 q^{k(t-k)}\le\gb tk\le2q^{k(t-k)},\qquad
 G_{t,k}\le2q^{-k(k+1)/2}.
\end{equation}
For the upper bound, divide the Gaussian coefficient by its leading power.
The numerator is at most one, whereas the denominator is at least
$\prod_{j\ge1}(1-q^{-j})\ge1-\sum_{j\ge1}q^{-j}\ge1/2$.
The lower bound holds term by term in the Gaussian product.

By \eqref{eq:mu-product}, $B_{a,k}$ vanishes unless $a-k$ is even;
otherwise $|B_{a,k}|^2\le q^{-(a-k)}$. Dropping the nonnegative subtracted
mean square in \eqref{eq:KD} gives
\[
 K_\D(r,t)\le2q^{-a}
 \sum_{\substack{k\ge0\\ k\equiv a\pmod2}}q^{-k(k-1)/2}
 \le4q^{-a}.
\]
Each parity-restricted series is at most two for $q\ge3$.

For $a=0$, $K_\J=p_r(1-p_r)\le9/(8q)$, which is enough.
Suppose $a\ge1$. Let $h(C)=\E[\J(Q|_U)\mid C]$ and compare it with
$p_a$, rather than its unknown-at-this-stage mean. If the common block
has nullity zero, $h(C)=p_a$. For $1\le k\le a$, product bounds give
\[
 |A_{a,k}-p_a|
 \le (1-\rho_{k,a})+(p_{a-k}-p_a)
 \le3q^{-(a-k+1)}.
\]
For $k>a$ the difference is at most one. Consequently
\begin{align*}
 K_\J(r,t)&=\Var h(C)\le\E|h(C)-p_a|^2\\
 &\le18q^{-2a-2}\sum_{k\ge1}q^{k(3-k)/2}
       +2\sum_{k\ge a+1}q^{-k(k+1)/2}\\
 &\le54q^{-2a-1}+3q^{-(a+1)(a+2)/2}
 \le57q^{-2a-1}\le64q^{-2a-1}.
\end{align*}
We used $\sum_{k\ge1}q^{k(3-k)/2}\le3q$,
and $(a+1)(a+2)/2\ge2a+1$ for $a\ge1$.

Finally fix $r,t\ge1$. If $a$ is even, the $k=0$ term in the second
moment in \eqref{eq:KD} equals
$p_t\mu_a^2=q^{-a}(1+O_r(q^{-1}))$; the remaining nonzero terms have
$k\ge2$ and are $O_r(q^{-a-1})$ in total. If $a$ is odd, the first
nonzero term is $k=1$, and
\[
 G_{t,1}p_{t-1}\rho_{1,a}^2\mu_{a-1}^2
 =q^{-a}(1+O_r(q^{-1})).
\]
The remaining nonzero terms have $k\ge3$ and are smaller. In either case
$\mu_r^2=O_r(q^{-r})=O_r(q^{-a-1})$ when nonzero.
The same argument includes $a=0$, where the leading term is $p_r$.
\end{proof}

\begin{theorem}[Exact profile formula for arbitrary subspace sampling]
\label{thm:profile}
Let $\mathcal U$ be any finite family of $r$-subspaces and let real weights
$\alpha_U$ have sum one. Put
$Z_f=\sum_{U\in\mathcal U}\alpha_U f(Q|_U)$.
Then $\E Z_\J=p_r$, $\E Z_\D=\mu_r$, $\E Z_{\I_\xi}=\pi_{r,\xi}$ and
\begin{equation}\label{eq:profile}
 \Var Z_f=\sum_{U,V\in\mathcal U}\alpha_U\alpha_V
 K_f\bigl(r,\dim(U\cap V)\bigr),
 \qquad f\in\{\J,\D,\I_\xi\}.
\end{equation}
The mixed covariance uses $K_{\J\D}$ in the same way.
For nonnegative weights the kernel upper bounds may be inserted term by
term. In particular,
\begin{equation}\label{eq:weighted-D-bound}
 \Var Z_\D\le4\sum_{\substack{U,V\in\mathcal U\\\dim(U\cap V)\ge1}}
 \alpha_U\alpha_V q^{-(r-\dim(U\cap V))}.
\end{equation}
\end{theorem}
\begin{proof}
Expand the variance or covariance of the finite weighted sum and apply
Corollary~\ref{cor:kernels}. All restrictions have the same marginal law.
For the last statement use nonnegative products of weights, the vanishing
at intersection dimension zero, and Lemma~\ref{lem:kernel-bounds}.
\end{proof}

This theorem retains the geometry suppressed by a uniform $O(q^{-1})$
off-diagonal estimate. It also handles repeated subspaces by combining
their weights. Signed weights are allowed in the exact identity, but not
in the termwise positive upper bound. In Section~\ref{sec:fourier} signed
weights instead lead to cancellation.

\section{Hyperbolic geometry and sharp all-generator averaging}
\label{sec:polar}
Let $V$ have dimension $2r$, and let $H$ be split and nondegenerate.
Thus $V$ has a hyperbolic basis $e_1,\ldots,e_r,f_1,\ldots,f_r$ in which
$H(e_i,e_j)=H(f_i,f_j)=0$ and $H(e_i,f_j)=\delta_{ij}$.
A \emph{generator} is a totally $H$-isotropic $r$-subspace.

\begin{lemma}[Intersection counts and the two families; classical]
\label{lem:polar-geometry}
The number of generators is
\begin{equation}\label{eq:N}
 N_r=2\prod_{j=1}^{r-1}(q^j+1).
\end{equation}
For a fixed generator $U$, the number of generators $W$ with
$\dim(U\cap W)=t$ is
\begin{equation}\label{eq:polar-count}
 n_{r,t}=\gb rt q^{(r-t)(r-t-1)/2}.
\end{equation}
There is a partition $\G=\G_+\sqcup\G_-$ into equal families, such that
$U,W$ are in the same family exactly when $r-\dim(U\cap W)$ is even.
Every proper totally isotropic subspace lies in the same number of
generators of each family.
\end{lemma}
\begin{proof}
Choose $T=U\cap W$ in $\gb rt$ ways. The quotient $T^\perp/T$ is a
split space of dimension $2(r-t)$, and $W/T$ is a generator complementary
to $U/T$. In a fixed complementary hyperbolic decomposition, such
subspaces are graphs of alternating matrices of size $r-t$: the graph
is isotropic precisely when the matrix plus its transpose is zero.
There are $q^{(r-t)(r-t-1)/2}$ choices, proving \eqref{eq:polar-count}.
Summing over $t$ and using
\[
 \prod_{j=0}^{r-1}(1+zq^j)
 =\sum_{a=0}^r q^{a(a-1)/2}\gb ra z^a
\]
at $z=1$ proves \eqref{eq:N}. The product identity follows by induction
from the Gaussian Pascal recurrence.

For completeness, the family partition can be constructed by determinants
of isometries. Fix a generator $U_0$. Every generator extends to a
hyperbolic basis: choose an isotropic basis vector $e$, choose a vector
$f$ with $H(e,f)=1$, replace $f$ by $f-H(f,f)e/2$, split off this
hyperbolic plane, and continue on its orthogonal complement. Thus an
isometry $g$ carries $U_0$ to any given $U$. In a hyperbolic decomposition
adapted to $U_0$, an isometry stabilizing $U_0$ has diagonal blocks
$A$ and $A^{-\mathsf T}$ and hence determinant one. Therefore
$s(U)=\det g\in\{1,-1\}$ is well defined.

Given two generators intersecting in dimension $t$, complete a basis of
the intersection to a hyperbolic basis in which the remaining $r-t$
directions of the two generators are paired complementary directions.
Swapping these $r-t$ hyperbolic pairs maps one generator to the other
and has determinant $(-1)^{r-t}$. This proves the parity criterion.
A single pair swap has determinant $-1$ and exchanges the two families,
so they have equal size.

If $T$ is proper and totally isotropic, choose a hyperbolic basis beginning
with a basis of $T$. A pair swap outside those initial coordinates fixes
$T$ pointwise and exchanges the families. It gives a bijection between
the two sets of generators containing $T$. This proves the final assertion.
\end{proof}

Set $w_{r,t}=n_{r,t}/N_r$. Equations~\eqref{eq:gaussian-bounds} and
\eqref{eq:N} imply the useful uniform estimate
\begin{equation}\label{eq:polar-weights}
 0\le w_{r,t}\le q^{-t(t-1)/2},\qquad \sum_{t=0}^r w_{r,t}=1.
\end{equation}

\begin{theorem}[Sharp polar averaging]\label{thm:polar}
Let $r\ge2$. With averages as in \eqref{eq:averages-intro}, their exact
variances are
\begin{equation}\label{eq:polar-exact}
 \Var(\overline f)=\sum_{t=0}^r w_{r,t}K_f(r,t),
 \qquad f\in\{\J,\D,\I_\xi\}.
\end{equation}
Uniformly in $r$ and odd $q$,
\begin{align}
 \Var(\overline\D)&\le12q^{-(r-1)},\label{eq:polar-D-bound}\\
 \Var(\overline\J)&\le192q^{-2(r-1)},\label{eq:polar-J-bound}\\
 \Var(\overline\I_\xi)&\le40q^{-(r-1)}.\label{eq:polar-I-bound}
\end{align}
For fixed $r$ as $q\to\infty$ through odd prime powers,
\begin{equation}\label{eq:polar-sharp}
 \Var(\overline\D)=(1+O_r(q^{-1}))q^{-(r-1)},\qquad
 \Var(\overline\I_\xi)=(\tfrac14+o_r(1))q^{-(r-1)}.
\end{equation}
In particular,
\begin{equation}\label{eq:polar-L1}
 \E|\overline\I_\xi-\pi_{r,\xi}|
 \le\sqrt{40}\,q^{-(r-1)/2},\qquad
 \E|\overline\I_\xi-\tfrac12|
 \le\sqrt{40}\,q^{-(r-1)/2}+q^{-1}.
\end{equation}
\end{theorem}
\begin{proof}
Apply Theorem~\ref{thm:profile} with uniform weights and the intersection
counts in Lemma~\ref{lem:polar-geometry}. The $t=0$ term vanishes.
For determinant character,
\[
 \Var(\overline\D)
 \le4q^{-(r-1)}\sum_{t\ge1}q^{-(t-1)(t-2)/2}
 \le12q^{-(r-1)}.
\]
For nonsingularity,
\[
 \Var(\overline\J)
 \le64q^{-2(r-1)}\sum_{t\ge1}q^{-(t-2)(t-3)/2}
 \le192q^{-2(r-1)}.
\]
Both displayed series are at most three for $q\ge3$.
Since $\overline\I_\xi=(\overline\J+\xi\overline\D)/2$,
\[
 \Var(\overline\I_\xi)
 \le\tfrac12\bigl(\Var(\overline\J)+\Var(\overline\D)\bigr)
 \le\bigl(6+96q^{-(r-1)}\bigr)q^{-(r-1)}
 \le40q^{-(r-1)}.
\]

For fixed $r$, the two leading contributions to the determinant variance
are $t=1$ and $t=2$. Indeed,
\[
 w_{r,1}=\tfrac12(1+O_r(q^{-1})),\qquad
 w_{r,2}=\tfrac12q^{-1}(1+O_r(q^{-1})),
\]
and Lemma~\ref{lem:kernel-bounds} gives respectively
$K_\D(r,1)=q^{-(r-1)}(1+O_r(q^{-1}))$ and
$K_\D(r,2)=q^{-(r-2)}(1+O_r(q^{-1}))$.
These formulas include the diagonal term when $r=2$.
The contributions from $t\ge3$ are $O_r(q^{-r})$.
Thus the two half-contributions add to the first asymptotic in
\eqref{eq:polar-sharp}. The nonsingularity variance is of smaller order,
and its covariance with $\overline\D$ is bounded by the geometric mean
of the variances. This proves the indicator asymptotic. The $L^1$ bounds
follow by Cauchy--Schwarz and Proposition~\ref{prop:marginals}.
\end{proof}

\begin{remark}[What is and is not improved]
For $r\ge3$, \eqref{eq:polar-L1} is $O(q^{-1})$, improving
\eqref{eq:source-bound} in the polar setting. For $r=2$, the true-centered
standard deviation remains of order $q^{-1/2}$. For fixed $r\ge3$,
$|\pi_{r,\xi}-1/2|=(2q)^{-1}+O_r(q^{-2})$, so an $o(q^{-1})$ bound
around $1/2$ is impossible without changing the center. The theorem does
not claim such a bound. When $r=1$, $N_1=2$ and ordinary polar averaging
has no concentration as $q\to\infty$; this is why $r\ge2$ is explicit.
\end{remark}

\section{Grassmannian averaging with an ambient-dimension exponent}
\label{sec:grassmann}
The same kernels give a second, quantitatively different law. Let
$\mathcal U=\Gr(r,\F_q^n)$, where $n\ge2r\ge2$, and let
$\overline f^{\,\mathrm{Gr}}$ denote the uniform average over this family.

\begin{theorem}[Sharp Grassmannian character averaging]\label{thm:grassmann}
For a fixed $r$-subspace, the fraction of $r$-subspaces meeting it in
dimension $t$ is
\begin{equation}\label{eq:grassmann-weights}
 v_{n,r,t}
 =\frac{\gb rt\gb{n-r}{r-t}q^{(r-t)^2}}{\gb nr}.
\end{equation}
The exact variance of each restriction statistic is
$\sum_{t=0}^r v_{n,r,t}K_f(r,t)$. Uniformly for $n\ge2r\ge2$,
\begin{align}
 \Var(\overline\D^{\,\mathrm{Gr}})&\le20q^{-(n-r)},\label{eq:grassmann-D}\\
 \Var(\overline\J^{\,\mathrm{Gr}})&\le384q^{-n},\label{eq:grassmann-J}\\
 \Var(\overline\I_\xi^{\,\mathrm{Gr}})&\le80q^{-(n-r)}.
 \label{eq:grassmann-I}
\end{align}
For fixed $n,r$ as $q\to\infty$,
\begin{equation}\label{eq:grassmann-asymptotic}
 \Var(\overline\D^{\,\mathrm{Gr}})
 =(1+O_{n,r}(q^{-1}))q^{-(n-r)},\quad
 \Var(\overline\I_\xi^{\,\mathrm{Gr}})
 =(\tfrac14+o_{n,r}(1))q^{-(n-r)}.
\end{equation}
For $n=2r$ and $r\ge2$,
$\E|\overline\I_\xi^{\,\mathrm{Gr}}-1/2|\le(\sqrt{80}+1)q^{-1}$.
\end{theorem}
\begin{proof}
Fix $U$. Choose $T=U\cap W$ in $\gb rt$ ways and the image of $W$
in $\F_q^n/U$ in $\gb{n-r}{r-t}$ ways. After choosing complements,
$W/T$ is the graph of a map from this image to $U/T$, giving
$q^{(r-t)^2}$ choices. This proves \eqref{eq:grassmann-weights}.
The variance identity follows from Theorem~\ref{thm:profile}.

The Gaussian bounds imply
$v_{n,r,t}\le4q^{-t(n-2r+t)}$.
Put $d=n-2r\ge0$. For determinant character, Lemma~\ref{lem:kernel-bounds}
therefore yields
\[
 \Var(\overline\D^{\,\mathrm{Gr}})
 \le16q^{-(n-r)}\sum_{t\ge1}q^{-(t-1)d-t(t-1)}
 \le20q^{-(n-r)}.
\]
For nonsingularity it yields
\[
 \Var(\overline\J^{\,\mathrm{Gr}})
 \le256q^{-n}\sum_{t\ge1}q^{-(t-1)d-(t-1)^2}
 \le384q^{-n}.
\]
The two series are at most $5/4$ and $3/2$, respectively, for $q\ge3$.
The indicator variance is at most half the sum of the other two, hence
at most $(10+192q^{-r})q^{-(n-r)}\le80q^{-(n-r)}$.

For the asymptotic, the unique leading term is $t=1$:
$v_{n,r,1}=q^{-(n-2r+1)}(1+O_{n,r}(q^{-1}))$ and
$K_\D(r,1)=q^{-(r-1)}(1+O_r(q^{-1}))$.
Terms with $t\ge2$ are smaller by at least $q^{-2}$ relative to this
leading scale when bounded as above. The nonsingularity variance is
smaller by a factor $q^{-r}$, so Cauchy--Schwarz proves the indicator
asymptotic. The final $L^1$ bound follows from the exact mean and
Proposition~\ref{prop:marginals}.
\end{proof}

\section{From incidence cancellation to Fourier cancellation}
\label{sec:fourier}
The sharp ordinary averages arose from positive overlap contributions.
We now use the difference of the two generator families to annihilate
those contributions at the Fourier level. This argument applies to
\emph{every} congruence-invariant statistic, not only determinant signs.

Fix a nontrivial additive character $\psi:\F_q\to\C^\times$. For
$f:\Sym_r\to\C$ define the normalized transform
\begin{equation}\label{eq:fourier-normalization}
 \widehat f(A)=q^{-c(r)}\sum_{C\in\Sym_r}f(C)\psi(\tr(AC)).
\end{equation}
The trace pairing is nondegenerate on symmetric matrices because $q$ is
odd: diagonal entries pair with coefficient one and off-diagonal entries
with coefficient two. Consequently
\begin{equation}\label{eq:fourier-parseval}
 f(C)=\sum_{A\in\Sym_r}\widehat f(A)\psi(-\tr(AC)),\qquad
 \E|f(C)|^2=\sum_{A\in\Sym_r}|\widehat f(A)|^2.
\end{equation}
A function is congruence-invariant when
$f(g^{\mathsf T}Cg)=f(C)$ for every $g\in\GL_r(\F_q)$.
Its Fourier transform is invariant under the dual congruence action as
well, by a change of variables in \eqref{eq:fourier-normalization}.

\begin{theorem}[Incidence-to-Fourier transfer]\label{thm:transfer}
Let $\mathcal U$ be a family of $r$-subspaces of a finite-dimensional
space $V$, let $w_U\in\C$, and let $f:\Sym_r\to\C$ be
congruence-invariant. Define
\[
 Z_w(Q)=\sum_{U\in\mathcal U}w_U f(Q|_U).
\]
Identify an ambient Fourier frequency with a symmetric tensor
$S\in\Sym^2V$, or its symmetric matrix in a chosen ambient basis, and
let $T=\im S$. Its Fourier coefficient in $Z_w$ is zero if no member
of $\mathcal U$ contains $T$. Otherwise it is
\begin{equation}\label{eq:transfer}
 \widehat {Z_w}(S)=\widehat f(A_{S,U})
               \sum_{\substack{W\in\mathcal U\\T\subseteq W}}w_W,
\end{equation}
where $U$ is any member containing $T$, and $A_{S,U}$ is the coordinate
matrix of $S$ inside $\Sym^2U$. Its congruence class, viewed in size $r$,
is independent of that choice.
In particular, if all incidence sums on subspaces of dimension at most
$s$ vanish, then $Z_w$ has no Fourier frequencies of rank at most $s$.
\end{theorem}
\begin{proof}
Let $P$ be the basis-column matrix of $U$. Restriction sends $Q$ to
$P^{\mathsf T}QP$, and
\[
 \tr(A P^{\mathsf T}QP)=\tr(PAP^{\mathsf T}Q).
\]
Thus lifting a restriction statistic embeds its frequency $A$ as
$S=PAP^{\mathsf T}$. The map $A\mapsto PAP^{\mathsf T}$ is injective;
its image consists exactly of symmetric tensors supported in $U$.
For example, a left inverse $R$ of $P$ gives
$A=RSR^{\mathsf T}$ whenever $\im S\subseteq U$.
In a basis of $U$ beginning with a basis of $T=\im S$, $A$ has a
nonsingular block supported on $T$ and a zero complementary block.
The same description holds in every other $r$-subspace containing $T$.
Changes of basis give congruent matrices, so congruence-invariance
makes $\widehat f(A)$ common to all these lifts. Summing the individual
Fourier coefficients proves \eqref{eq:transfer} and the annihilation
assertion. This includes the zero frequency and the zero subspace.
\end{proof}

\begin{theorem}[Exact top-rank energy identity]\label{thm:top-rank}
Let $H$ be split on $\F_q^{2r}$, $r\ge1$, and let $f:\Sym_r\to\C$ be
congruence-invariant. Define its two-family contrast by
\[
 \Delta_f(Q)=\frac2{N_r}\sum_{U\in\G_+}f(Q|_U)
             -\frac2{N_r}\sum_{U\in\G_-}f(Q|_U).
\]
Every Fourier coefficient of $\Delta_f$ of rank less than $r$ vanishes.
All its remaining frequencies have rank exactly $r$, with image a
generator. Moreover
\begin{equation}\label{eq:top-energy}
 \boxed{\displaystyle
 \E|\Delta_f|^2=\frac4{N_r}
             \sum_{\substack{A\in\Sym_r\\\det A\ne0}}|\widehat f(A)|^2.}
\end{equation}
For a second invariant statistic $g$,
\begin{equation}\label{eq:top-cross-energy}
 \E[\Delta_f\overline{\Delta_g}]
 =\frac4{N_r}\sum_{\det A\ne0}\widehat f(A)\overline{\widehat g(A)}.
\end{equation}
\end{theorem}
\begin{proof}
Use Theorem~\ref{thm:transfer} with weights $2/N_r$ and $-2/N_r$ on the
two families. If a frequency has rank less than $r$ and is supported on
some generator, its support is a proper totally isotropic subspace.
Lemma~\ref{lem:polar-geometry} makes the signed incidence sum zero.
If its rank is $r$, its image itself must be the unique generator
supporting it. Its coefficient is then $\pm(2/N_r)\widehat f(A)$ with
$A$ nonsingular. Frequencies of rank greater than $r$ never arise from
an $r$-dimensional restriction.

For each generator there are all nonsingular $A\in\Sym_r$, and these
full-rank frequency sets are disjoint for distinct generators, since
their images differ. Parseval therefore gives
$N_r(2/N_r)^2\sum_{\det A\ne0}|\widehat f(A)|^2$, proving
\eqref{eq:top-energy}. The mixed identity follows by the same orthogonality
calculation. All contrasts have mean zero since their zero frequency
vanishes.
\end{proof}

\begin{remark}[Why class invariance matters]
For arbitrary functions chosen in unrelated bases on the different
subspaces, the coefficient $\widehat f(A_{S,U})$ need not be constant
among the subspaces containing $\im S$. Equality of incidence counts
alone would then not prove cancellation. Congruence-invariance is an
essential hypothesis of the transfer theorem as stated.
\end{remark}

\begin{theorem}[Rigidity of the universal top-rank filter]
\label{thm:filter-rigidity}
Let weights $w_U$ be assigned to all generators of a split $2r$-space.
The following are equivalent: for every congruence-invariant $f$, the
weighted statistic $\sum_Uw_Uf(Q|_U)$ has no Fourier frequencies of rank
less than $r$; and $w_U=a$ on $\G_+$ and $w_U=-a$ on $\G_-$ for one scalar
$a$. Thus the universal maximal-rank filter in this generator family is
unique up to scaling.
\end{theorem}
\begin{proof}
The indicated weights have the property by Theorem~\ref{thm:top-rank}.
Conversely, test the asserted vanishing with $f=\delta_0$, the indicator
of the zero form, whose Fourier coefficients are all $q^{-c(r)}$.
For each totally isotropic $(r-1)$-space $T$, choose a frequency with
image $T$. Equation~\eqref{eq:transfer} forces the sum of weights on
generators containing $T$ to vanish. There are exactly two such
generators, from the split rank-one quotient $T^\perp/T$, so the weights
are opposite on adjacent generators. The adjacency graph is connected:
in a hyperbolic basis adapted to two generators, swap their differing
paired directions one at a time, obtaining a path of codimension-one
intersections. Propagating the opposite-weight relation along such a
path and using the parity criterion in Lemma~\ref{lem:polar-geometry}
proves the claim. For $r=1$, the same argument uses the zero subspace
and the two generators directly.
\end{proof}

\section{Exact two-family discrepancies}
\label{sec:contrast}
\subsection{The nonsingular Fourier coefficients}
We evaluate only the full-rank Fourier layer needed by
Theorem~\ref{thm:top-rank}. These evaluations lie within the classical
symmetric-form Fourier theory \cite{schmidt2015,schmidt2020}; the elementary
recursion is included to fix every normalization and parity sign.
Put
\[
 \gamma=\sum_{x\in\F_q}\chi(x)\psi(x).
\]
The usual one-variable calculation gives
\begin{equation}\label{eq:gauss}
 |\gamma|^2=q,\qquad \gamma^2=\eps q,\qquad
 \sum_x\psi(ax^2)=\chi(a)\gamma\quad(a\ne0).
\end{equation}
For clarity, the first identity follows by substituting $x=uy$ in
$\sum_{x,y}\chi(x)\chi(y)\psi(x-y)$: the inner sum over $y\ne0$ is
$q-1$ for $u=1$ and $-1$ otherwise. Also
$\overline\gamma=\eps\gamma$, which gives the second identity.
The last identity follows by counting the two square roots of each
nonzero square and using $\sum_x\psi(ax)=0$.

Define unnormalized transforms
\[
 d_n(A)=\sum_{C\in\Sym_n}\D(C)\psi(\tr(AC)),\qquad
 j_n(A)=\sum_{C\in\Sym_n}\J(C)\psi(\tr(AC)).
\]
The empty values are $d_0=j_0=1$.

\begin{lemma}[Full-rank Fourier recursion]\label{lem:fourier-recursion}
For a nonsingular block diagonal matrix $A=B\oplus a$, $a\ne0$,
\begin{equation}\label{eq:fourier-recursions}
 d_n(B\oplus a)=\chi(a)^n\gamma^n j_{n-1}(B),\qquad
 j_n(B\oplus a)=-\chi(a)^{n-1}\gamma^{n-1}d_{n-1}(B).
\end{equation}
For nonsingular $A$ these imply
\begin{align}
 d_{2m}(A)&=(-1)^m\gamma^{2m^2},&
 j_{2m}(A)&=\chi(\det A)d_{2m}(A),\label{eq:fourier-even}\\
 d_{2m+1}(A)&=(-1)^m\chi(\det A)\gamma^{2m^2+2m+1},&
 j_{2m+1}(A)&=(-1)^{m+1}\gamma^{2m^2+2m}.
 \label{eq:fourier-odd}
\end{align}
Consequently, at every nonsingular frequency,
\begin{equation}\label{eq:fourier-magnitudes}
 |\widehat\D_n(A)|^2=q^{-c(n)-\lfloor n/2\rfloor},\qquad
 |\widehat\J_n(A)|^2=q^{-c(n)-\lceil n/2\rceil}.
\end{equation}
\end{lemma}
\begin{proof}
Write
$C=\left(\begin{smallmatrix}C_0&z\\z^{\mathsf T}&x\end{smallmatrix}\right)$.
If $C_0$ is singular, $\det C$ is independent of $x$, since its coefficient
as a polynomial in $x$ is $\det C_0=0$. Summing $\psi(ax)$ then makes
its entire contribution to both transforms zero.

If $C_0$ is nonsingular, set $h=z^{\mathsf T}C_0^{-1}z$. Then
$\det C=\det C_0(x-h)$. The sum over $x$ in the determinant-character
transform is
\[
 \chi(\det C_0)\chi(a)\gamma\,\psi(ah),
\]
whereas the sum in the nonsingularity transform is $-\psi(ah)$.
Diagonalizing $C_0^{-1}$ and using \eqref{eq:gauss} gives
\[
 \sum_z\psi(a z^{\mathsf T}C_0^{-1}z)
 =\chi(a)^{n-1}\chi(\det C_0)\gamma^{n-1}.
\]
Multiplying the factors and summing over nonsingular $C_0$ with weight
$\psi(\tr(BC_0))$ proves \eqref{eq:fourier-recursions}.
Every nonsingular symmetric matrix is congruent to a nonsingular diagonal
matrix: choose a vector of nonzero norm and split off its line repeatedly,
using polarization in odd characteristic. Thus congruence-invariance
allows the recursion to be iterated for every nonsingular $A$.
The formulas \eqref{eq:fourier-even}--\eqref{eq:fourier-odd} follow by
induction, starting at $n=0$. Squaring magnitudes, using
$|\gamma|^2=q$, and dividing by $q^{2c(n)}$ proves
\eqref{eq:fourier-magnitudes}.
\end{proof}

\subsection{The exact second moments}
\begin{theorem}[Two-family discriminant and sign discrepancy]\label{thm:contrast}
Let $r\ge1$, $q$ odd, and let $\Delta_f$ be the contrast between the two
families of generators of a fixed split form on $\F_q^{2r}$. Then
\begin{align}
 \E\Delta_\D^2&=\frac{4p_r}{N_r}q^{-\lfloor r/2\rfloor},
 \label{eq:contrast-D}\\
 \E\Delta_\J^2&=\frac{4p_r}{N_r}q^{-\lceil r/2\rceil},
 \label{eq:contrast-J}\\
 \E[\Delta_\J\Delta_\D]
 &=\begin{cases}
 \displaystyle\frac{4\mu_r}{N_r}q^{-r/2},&r\text{ even},\\[4pt]
 0,&r\text{ odd}.
 \end{cases}\label{eq:contrast-cross}
\end{align}
The individual-sign contrasts satisfy
\begin{equation}\label{eq:contrast-I}
 \E\Delta_{\I_\xi}^2
 =\begin{cases}
 \displaystyle\frac{2q^{-m}}{N_r}(p_r+\xi\mu_r)
       =\frac{4q^{-m}}{N_r}\pi_{r,\xi},&r=2m,\\[6pt]
 \displaystyle\frac{p_rq^{-m}}{N_r}(1+q^{-1}),&r=2m+1.
 \end{cases}
\end{equation}
All these contrasts have mean zero, so these second moments are variances.
\end{theorem}
\begin{proof}
There are $q^{c(r)}p_r$ nonsingular matrices in $\Sym_r$.
Insert \eqref{eq:fourier-magnitudes} in \eqref{eq:top-energy} to obtain
\eqref{eq:contrast-D} and \eqref{eq:contrast-J}.
For even $r$, \eqref{eq:fourier-even} gives
$\widehat\J(A)=\chi(\det A)\widehat\D(A)$ on the nonsingular layer.
The sum of $\chi(\det A)$ on that layer is $q^{c(r)}\mu_r$, so
\eqref{eq:top-cross-energy} gives the first case of
\eqref{eq:contrast-cross}. For odd $r$, the mixed Fourier product is a
constant multiple of $\chi(\det A)$, whose sum is zero since $\mu_r=0$.
Finally $\Delta_{\I_\xi}=(\Delta_\J+\xi\Delta_\D)/2$; expanding its
square gives \eqref{eq:contrast-I}.
\end{proof}

\begin{corollary}[Quadratic exponent in both asymptotic regimes]
\label{cor:quadratic}
For every odd $q$ and $r\ge1$,
\begin{equation}\label{eq:normalized-contrast}
 q^{\lfloor r^2/2\rfloor}\E\Delta_\D^2
 =\frac{2p_r}{\prod_{j=1}^{r-1}(1+q^{-j})}\le2.
\end{equation}
For fixed $r$ this quantity is $2+O_r(q^{-1})$ as $q\to\infty$.
For fixed odd $q$, let
$p_\infty(q)=\prod_{j\ge1,\ j\text{ odd}}(1-q^{-j})>0$.
Then
\begin{equation}\label{eq:fixed-q-limit}
 \lim_{r\to\infty}q^{\lfloor r^2/2\rfloor}\E\Delta_\D^2
 =2p_\infty(q)^2.
\end{equation}
\end{corollary}
\begin{proof}
Use $N_r=2q^{r(r-1)/2}\prod_{j=1}^{r-1}(1+q^{-j})$ and the integer
identity $r(r-1)/2+\lfloor r/2\rfloor=\lfloor r^2/2\rfloor$.
This gives \eqref{eq:normalized-contrast} and the fixed-$r$ limit.
The products converge absolutely for fixed $q$. Writing $x=q^{-1}$,
\[
 \prod_{j\ge1}(1+x^j)
 =\frac{\prod_{j\ge1}(1-x^{2j})}{\prod_{j\ge1}(1-x^j)}
 =\frac1{\prod_{j\ge1,\ j\text{ odd}}(1-x^j)}.
\]
Taking the limit in \eqref{eq:normalized-contrast} therefore proves
\eqref{eq:fixed-q-limit}.
\end{proof}

\begin{corollary}[Common fluctuations and almost perfect correlation]
\label{cor:correlation}
For $f=\J,\D$, or $\I_\xi$,
\begin{equation}\label{eq:family-variance}
 \Cov(\overline f,\Delta_f)=0,\qquad
 \Var(F_{f,\pm})=\Var(\overline f)+\tfrac14\E\Delta_f^2.
\end{equation}
For determinant character, fixed $r\ge2$, and $q\to\infty$,
\begin{equation}\label{eq:family-correlation}
 1-\operatorname{Corr}(F_{\D,+},F_{\D,-})
 =(1+O_r(q^{-1}))q^{-\lfloor r^2/2\rfloor+r-1}.
\end{equation}
\end{corollary}
\begin{proof}
An $H$-isometry of determinant $-1$ exchanges the families and preserves
the uniform law of $Q$ under congruence. It fixes $\overline f$ and
negates $\Delta_f$, so their covariance is zero.
Use $F_{f,\pm}=\overline f\pm\Delta_f/2$ to obtain
\eqref{eq:family-variance} and
$\Cov(F_{f,+},F_{f,-})=\Var(\overline f)-\E\Delta_f^2/4$.
For $f=\D$, divide by the common variance of the two family averages.
If $v=\Var(\overline\D)$ and $d=\E\Delta_\D^2$, the correlation deficit
is $d/(2v+d/2)$. Equations~\eqref{eq:polar-sharp} and
\eqref{eq:normalized-contrast} give \eqref{eq:family-correlation}.
\end{proof}

For example, Markov's inequality applied to the exact second moment gives
\begin{equation}\label{eq:contrast-tail}
 \Pp(|\Delta_\D|\ge u)
 \le\min\left\{1,\frac{4p_rq^{-\lfloor r/2\rfloor}}{N_ru^2}\right\}
 \quad(u>0).
\end{equation}
This is a rigorous second-moment tail bound, not a claim of Gaussian tails
or a central limit theorem. The family means remain
$\E F_{\D,\pm}=\mu_r$, $\E F_{\J,\pm}=p_r$, and
$\E F_{\I_\xi,\pm}=\pi_{r,\xi}$.

\section{Rare geometric exceptions and a non-Gaussian limit}
\label{sec:rare}
The exact variance has a stronger pointwise explanation in one parity for
each observable. In even generator dimension, the determinant-character
averages agree \emph{exactly} unless the two forms have a common generator.
In odd generator dimension the same statement holds for nonsingularity.
This reveals why a standard deviation need not describe a typical sample.

For a fixed split $H$ and arbitrary symmetric $Q$, define
\begin{equation}\label{eq:common-counts}
 C_\pm(Q)=|\{U\in\G_\pm:Q|_U=0\}|,\qquad
 C(Q)=C_+(Q)+C_-(Q),\qquad Z(Q)=C_+(Q)-C_-(Q).
\end{equation}
Here $Q|_U=0$ means that the whole restricted form is zero, not merely
that its determinant vanishes. Set $M_r=N_r/2$.

\begin{theorem}[Pointwise geometric discrepancy identity]
\label{thm:pointwise}
For every symmetric $Q$,
\begin{align}
 \Delta_\D(Q)&=(-\eps)^m\frac{q^{m^2}}{M_r}Z(Q),
 &&r=2m,\label{eq:pointwise-D}\\
 \Delta_\J(Q)&=(-1)^{m+1}\frac{q^{m(m+1)}}{M_r}Z(Q),
 &&r=2m+1.\label{eq:pointwise-J}
\end{align}
Consequently, the corresponding contrast is zero whenever $H$ and $Q$
have no common totally isotropic $r$-subspace.
If either displayed contrast is nonzero, the characteristic polynomial
of the $H$-self-adjoint operator $T=H^{-1}Q$ is a square in $\F_q[x]$.
In particular, a nonsquare characteristic polynomial certifies exact
equality of the corresponding two family averages.
\end{theorem}
\begin{proof}
Let $\delta_0$ be the indicator of the zero matrix on $\Sym_r$.
Its Fourier transform is identically $q^{-c(r)}$, and its contrast is
$\Delta_{\delta_0}=Z/M_r$.
For even $r=2m$, Lemma~\ref{lem:fourier-recursion} shows that the
unnormalized transform of $\D$ on every nonsingular frequency equals
\[
 (-1)^m\gamma^{2m^2}=(-1)^m\eps^{m^2}q^{m^2}
 =(-\eps)^m q^{m^2}.
\]
Theorem~\ref{thm:top-rank} annihilates all other Fourier frequencies of
the contrast. Thus the transforms of $\Delta_\D$ and
$(-\eps)^m q^{m^2}\Delta_{\delta_0}$ agree everywhere. Fourier inversion
proves \eqref{eq:pointwise-D}. For odd $r=2m+1$, the full-rank
unnormalized transform of $\J$ is
\[
 (-1)^{m+1}\gamma^{2m^2+2m}=(-1)^{m+1}q^{m(m+1)},
\]
since $m(m+1)$ is even. The same argument proves
\eqref{eq:pointwise-J}.

A nonzero discrepancy forces $Z\ne0$, hence a common generator $U$.
Since $Q(u,v)=H(Tu,v)$, its vanishing on $U\times U$ implies
$TU\subseteq U^{\perp_H}=U$. In a hyperbolic basis adapted to $U$,
self-adjointness then forces the block form
\[
 T=\begin{pmatrix}A&B\\0&A^{\mathsf T}\end{pmatrix},\qquad B=B^{\mathsf T}.
\]
Therefore $\det(xI-T)=\det(xI-A)^2$.
The square condition does not conversely force a nonzero discrepancy:
for $Q=aH$, all generators are common and $C_+=C_-=M_r$.
\end{proof}

\begin{theorem}[Exact rare-event moments and asymptotic law]
\label{thm:rare-law}
For uniform $Q$, put
\begin{equation}\label{eq:lambda-R}
 \Lambda_r=N_rq^{-c(r)},\qquad
 R_r=\sum_{a=1}^r\gb ra q^{-a(r+1-a)}.
\end{equation}
Then
\begin{equation}\label{eq:common-moments}
 \E C=\Lambda_r,\qquad
 \E[C(C-1)]=\Lambda_r R_r,\qquad
 \E Z=0,\qquad \E Z^2=\Lambda_r p_r.
\end{equation}
In particular,
\begin{equation}\label{eq:rare-bound}
 \Pp(Z\ne0)\le\Lambda_r
 =2q^{-r}\prod_{j=1}^{r-1}(1+q^{-j})\le4q^{-r}.
\end{equation}
For fixed $r\ge1$ as $q\to\infty$,
\begin{align}
 \Pp(Z=1)=\Pp(Z=-1)&=(1+O_r(q^{-1}))q^{-r},
 \label{eq:rare-one}\\
 \Pp(Z\ne0)&=(2+O_r(q^{-1}))q^{-r},
 \label{eq:rare-nonzero}\\
 \Pp(|Z|\ge2)&=O_r(q^{-r-1}).\label{eq:rare-many}
\end{align}
Thus the selected contrasts in Theorem~\ref{thm:pointwise} are exactly
zero except on a rare geometric event, and their first-order nonzero
values are the two opposite multiples of one specified there.
\end{theorem}
\begin{proof}
For each generator $U$, restriction is uniform on $\Sym_r$, so
$\Pp(Q|_U=0)=q^{-c(r)}$. This proves $\E C=\Lambda_r$.
For two generators with intersection dimension $t$, the two zero
restrictions impose $2c(r)-c(t)$ independent linear conditions.
Using \eqref{eq:polar-count} and setting $a=r-t$ gives
\begin{align*}
 \E[C(C-1)]
 &=N_r\sum_{t=0}^{r-1}n_{r,t}q^{-2c(r)+c(t)}\\
 &=N_rq^{-c(r)}\sum_{a=1}^r\gb ra q^{-a(r+1-a)}
 =\Lambda_r R_r.
\end{align*}
A family-swapping isometry gives $\E Z=0$. To compute its second
moment, apply \eqref{eq:top-energy} to $\delta_0$:
\[
 \E\Delta_{\delta_0}^{\,2}
 =\frac4{N_r}\bigl(q^{c(r)}p_r\bigr)q^{-2c(r)}.
\]
Since $Z=M_r\Delta_{\delta_0}$, this equals the final identity in
\eqref{eq:common-moments}.

The event $Z\ne0$ is contained in $C\ge1$, so Markov's inequality gives
the first bound in \eqref{eq:rare-bound}.
The product is at most
$\exp(\sum_{j\ge1}q^{-j})\le e^{1/2}<2$, proving its uniform version.
For the asymptotic law, \eqref{eq:gaussian-bounds} gives
$R_r\le2\sum_{a\ge1}q^{-a}\le2/(q-1)$, and the $a=1$ term gives
$R_r=q^{-1}(1+O_r(q^{-1}))$.
For any nonnegative integer $C$,
\[
 \Pp(C\ge2)\le\tfrac12\E[C(C-1)],\qquad
 \E C-\E[C(C-1)]\le\Pp(C=1)\le\E C.
\]
It follows that
$\Pp(C=1)=\Lambda_r(1+O_r(q^{-1}))$ and
$\Pp(C\ge2)=O_r(\Lambda_r q^{-1})$.
When $C=1$, $Z=1$ or $-1$; the two cases have equal probability by
family-swapping symmetry. The difference between these single-common-
generator events and $Z=\pm1$ is contained in $C\ge2$.
Finally $\Lambda_r=(2+O_r(q^{-1}))q^{-r}$. This proves
\eqref{eq:rare-one}--\eqref{eq:rare-many}.
\end{proof}

\begin{corollary}[Variance normalization does not produce a Gaussian limit]
\label{cor:non-gaussian}
Let $f_r=\D$ when $r$ is even and $f_r=\J$ when $r$ is odd, and set
\[
 W_{r,q}=\frac{\Delta_{f_r}}{\sqrt{\E\Delta_{f_r}^{\,2}}}.
\]
Then $\E W_{r,q}=0$ and $\E W_{r,q}^2=1$ for every $r,q$, but
$W_{r,q}\to0$ in probability both when $r$ is fixed and $q\to\infty$,
and when $q$ is fixed and $r\to\infty$. In these regimes its limiting
distribution is the point mass at zero, not a nondegenerate Gaussian,
and its squares are not uniformly integrable.
\end{corollary}
\begin{proof}
By Theorem~\ref{thm:pointwise}, $W_{r,q}$ is, up to a sign,
$Z/\sqrt{\Lambda_r p_r}$. Thus it has second moment one by
\eqref{eq:common-moments}, while
$\Pp(W_{r,q}\ne0)\le\Lambda_r\le4q^{-r}$ tends to zero in both
stated regimes. This proves convergence in probability. Uniform
integrability of the squares would then force their expectations to
tend to zero, contradicting their constant value one.
\end{proof}

\begin{remark}[A substantive limit-law distinction]
For even $r$, the variance $\E\Delta_\D^2\asymp q^{-r^2/2}$ is carried
by jumps of order $q^{-r^2/4+r/2}$ occurring with probability of order
$q^{-r}$. The jump size is therefore much larger than the root mean
square by a factor of order $q^{r/2}$. The odd-rank determinant-character
contrast and the even-rank nonsingularity contrast are not covered by
this sparse pointwise identity; no Gaussian or non-Gaussian limit is
asserted for those remaining cases.
\end{remark}

\begin{theorem}[Complete rank-two discrepancy law]\label{thm:rank-two-law}
For $r=2$, put
\[
 a_q=\frac{(q^2-1)(q^3-q+1)}{q^7},\qquad
 b_q=\frac{(q-1)^2(q+1)}{4q^7}.
\]
The signed common-generator count $Z$ has exactly the following law:
\begin{center}
\begin{tabular}{@{}cccccc@{}}
\toprule
Value of $Z$&$-2$&$-1$&$0$&$1$&$2$\\
\midrule
Probability&$b_q$&$a_q$&$1-2a_q-2b_q$&$a_q$&$b_q$\\
\bottomrule
\end{tabular}
\end{center}
In particular $\Delta_\D=-\eps qZ/(q+1)$, so this also gives the full
finite-$q$ law of the determinant-character contrast, not merely its
second moment.
\end{theorem}
\begin{proof}
Choose coordinates in which the projective quadric of $H$ is
$\det X=0$ for a $2$-by-$2$ matrix $X=(x_{ij})$.
Its points have the form $X=uv^{\mathsf T}$ with $[u],[v]\in\mathbb P^1$.
Its two line families are the fibers with one factor fixed: indeed a
line spanned by $uv^{\mathsf T}$ and $u'v'^{\mathsf T}$ consists of
rank-one matrices only if
$\det(u,u')\det(v,v')=0$, forcing one factor to be fixed.
Pulling back the quadratic form $Q$ gives a bihomogeneous polynomial
$F(u,v)$ of bidegree $(2,2)$. The resulting linear map from the
10-dimensional space of quadratics in the four $x_{ij}$ to the
9-dimensional space of bidegree-$(2,2)$ polynomials is onto: every
bihomogeneous monomial is a product of two $x_{ij}=u_iv_j$.
Its kernel is therefore the one-dimensional span of $H$, and uniform
$Q$ induces uniform $F$ in this 9-dimensional space.

If $F=0$, both family counts are $q+1$ and $Z=0$. If $F\ne0$, at most
two fibers from either family can lie in its zero set, since a nonzero
binary quadratic coefficient has at most two projective roots.
Hence $|Z|\le2$. Having $Z=2$ means that exactly two fibers from the
first family and no fibers from the second lie in the zero set.
The first condition forces
$F(u,v)=a(u)b(v)$, where $a$ is a binary quadratic with two distinct
$\F_q$-projective roots; the second says that $b$ has no such root.

The number of choices of $a$ is
$(q-1)\binom{q+1}{2}=q(q^2-1)/2$.
The number of nonzero binary quadratics $b$ with no projective root is
$q(q-1)^2/2$, by Proposition~\ref{prop:marginals} in dimension two
(or by counting irreducible affine quadratics).
A nonzero tensor $a\otimes b$ has $q-1$ such factorizations, from
opposite scalar rescalings. Division by $q^9$ therefore gives
\[
 \Pp(Z=2)=\frac{q(q^2-1)}2\frac{q(q-1)^2}2
          \frac1{q-1}\frac1{q^9}
 =b_q.
\]
The opposite value has the same probability. By
\eqref{eq:common-moments}, $\E Z^2=2(q^2-1)/q^4$.
Symmetry and the five-point support imply
$2\Pp(Z=1)+8b_q=2(q^2-1)/q^4$, which gives $a_q$.
The remaining probability is its complement. The final contrast
identity is \eqref{eq:pointwise-D} with $m=1$.
\end{proof}

\section{A local improvement in the rank-layer transverse argument}
\label{sec:application}
This section makes the source interface explicit and proves it in the
present notation. It improves a local estimate, not the separate global
claims, singular-tail estimates, or parameter optimization of the cited
manuscript.

\subsection{Conditioning without an independence assumption}
\begin{lemma}[High-probability conditioning]\label{lem:conditioning}
Let $Q$ be uniform in $\Sym_n$ and $L$ fixed. The event
$\mathcal A_L=\{Q\text{ and }Q-L\text{ nonsingular}\}$ satisfies
\begin{equation}\label{eq:conditioning-prob}
 \Pp(\mathcal A_L)\ge1-2(1-p_n)\ge1-\frac9{4q}\ge\frac14.
\end{equation}
For every nonnegative random variable $Z$,
\begin{equation}\label{eq:conditioning-bound}
 \E[Z\mid\mathcal A_L]\le\frac{\E Z}{\Pp(\mathcal A_L)}\le4\E Z.
\end{equation}
The same conclusions hold when only $Q$ is required to be nonsingular.
In particular the $O(q^{-1})$ polar $L^1$ bound for $r\ge3$, the
$O(q^{-1})$ middle-Grassmannian bound for $r\ge2$, and the quadratic
contrast second-moment upper bound persist under these conditionings.
\end{lemma}
\begin{proof}
Both $Q$ and $Q-L$ are marginally uniform. Apply the union bound and
Proposition~\ref{prop:marginals}; no independence is required.
For nonnegative $Z$, $\E[Z\one_{\mathcal A_L}]\le\E Z$, proving
\eqref{eq:conditioning-bound}. Insert the appropriate absolute error
or contrast square to obtain the stated consequences. This argument
does not assert that the exact conditional mean or exact conditional
variance equals the unconditional one.
\end{proof}

\begin{lemma}[Relative conditioning along a nonsingular pencil]
\label{lem:pencil-conditioning}
Let $H\in\Sym_n$ be nonsingular, let $S\subseteq\F_q$ have size $k$,
and put
\[
 \mathcal A_S=\{Q-sH\text{ nonsingular for every }s\in S\},\qquad
 \delta=kn/q<1.
\]
If a nonnegative statistic $X(Q)$ is invariant under every translation
$Q\mapsto Q+aH$, then
\begin{equation}\label{eq:relative-conditioning}
 (1-\delta)\E X\le\E[X\mid\mathcal A_S]
                  \le(1-\delta)^{-1}\E X.
\end{equation}
For a real translation-invariant statistic $Y$, with $v=\Var Y$,
\begin{align}
 |\E[Y\mid\mathcal A_S]-\E Y|
 &\le\frac{\delta}{1-\delta}\sqrt v,\label{eq:relative-mean}\\
 \left(1-\delta-\frac{\delta^2}{(1-\delta)^2}\right)v
 \le\Var(Y\mid\mathcal A_S)&\le(1-\delta)^{-1}v.
 \label{eq:relative-variance}
\end{align}
\end{lemma}
\begin{proof}
Partition $\Sym_n$ into its affine lines $Q+\F_q H$.
On each line, $\det(Q+aH)$ is a polynomial in $a$ of degree $n$ with
nonzero leading coefficient $\det H$, so it has at most $n$ roots.
The $k$ shifted nonsingularity conditions exclude at most $kn$ values
of $a$. Thus, conditional on the affine line, the probability of
$\mathcal A_S$ lies between $1-\delta$ and one. A line-invariant
nonnegative $X$ can be multiplied by these bounds before averaging;
dividing by $\Pp(\mathcal A_S)\in[1-\delta,1]$ proves
\eqref{eq:relative-conditioning}.

Let $Y_0=Y-\E Y$. Its mean is zero, and subtracting one from the
conditional line probability shows
$|\E[Y_0\one_{\mathcal A_S}]|\le\delta\E|Y_0|\le\delta\sqrt v$.
This proves \eqref{eq:relative-mean}. Apply
\eqref{eq:relative-conditioning} to $Y_0^2$, then subtract the square
of the conditional mean to obtain \eqref{eq:relative-variance}.
\end{proof}

\begin{corollary}[The sharp polar laws survive pencil conditioning]
\label{cor:conditional-polar}
For a fixed split $H$ on $\F_q^{2r}$ and a fixed finite set of shifts
$S$, the expectations of the squared contrasts $\Delta_f^2$ for
$f\in\{\J,\D,\I_+,\I_-\}$ change from their exact values in
Theorem~\ref{thm:contrast} by a factor $1+O_{r,|S|}(q^{-1})$
after conditioning on $\mathcal A_S$. The conditional means of the
contrasts are zero.
The variances of the ordinary polar averages also change by a factor
$1+O_{r,|S|}(q^{-1})$.
For the sparse contrasts of Theorem~\ref{thm:pointwise}, the conditional
probability of a nonzero discrepancy is
$(2+O_{r,|S|}(q^{-1}))q^{-r}$.
\end{corollary}
\begin{proof}
Every restriction to an $H$-isotropic generator is unchanged by
$Q\mapsto Q+aH$. Apply Lemma~\ref{lem:pencil-conditioning} to contrast
squares, ordinary averages, and $\one_{\{Z\ne0\}}$.
An $H$-isometry exchanging the two families preserves
$\mathcal A_S$ and negates each contrast, so its conditional mean is
exactly zero. The probability assertion follows from
\eqref{eq:rare-nonzero}. The estimates concern fixed $r,|S|$ as
$q\to\infty$; the explicit lemma requires $q>2r|S|$.
No relative preservation of the mixed moment
$\E[\Delta_\J\Delta_\D]$ is asserted.
\end{proof}

\subsection{The local statistic and its improved estimate}
For $D=2r$ let $\mathcal R_{r,\xi}\subseteq\Sym_D$ consist of matrices of
rank $r$ and nondegenerate-quotient determinant character $\xi$. Write
\begin{equation}\label{eq:rank-layer-normalization}
 \lambda_{r,\xi}=\frac{|\mathcal R_{r,\xi}|}{q^{c(D)}}
 =\gb Dr q^{c(r)-c(D)}\pi_{r,\xi},\qquad
 K_\xi(A)=\frac{\one_{\mathcal R_{r,\xi}}(A)}{\lambda_{r,\xi}}.
\end{equation}
The count follows by choosing the image of a rank-$r$ symmetric matrix:
for a basis-column matrix $P$ of its image it has a unique representation
$PRP^{\mathsf T}$ with $R$ nonsingular symmetric, and its quotient sign
is $\chi(\det R)$.
For independent uniform $X_1,X_2,X_3\in\Sym_D$, signs
$\xi_1,\xi_2,\xi_3$, and an active subset $A\subseteq\{1,2,3\}$, define
\begin{align}
 F_A(X)&=\E_Y\prod_{i\in A}K_{\xi_i}(X_i-Y),\label{eq:local-F}\\
 G_A(X)&=\prod_{\substack{i<j\\i,j\in A}}
 \bigl(1+\eps^r\chi(\det(X_i-X_j))\xi_i\xi_j\bigr),
 \label{eq:local-G}\\
 \mathcal T&=\{X_i-X_j\text{ nonsingular for all }i<j\}.
 \label{eq:local-T}
\end{align}
The empty products are one. On $\mathcal T$, $G_A$ is zero on a failed
active pair match and otherwise equals $2^{\binom{|A|}{2}}$.

\begin{theorem}[Improved local transverse approximation]\label{thm:local}
For fixed $r\ge3$, uniformly in the active subset and signs,
\begin{equation}\label{eq:local-improvement}
 \E\bigl[\one_{\mathcal T}|F_A-G_A|\bigr]=O_r(q^{-1})
 \quad(q\to\infty\text{ through odd prime powers}).
\end{equation}
On $\mathcal T$ one also has $0\le F_A\le C_r$, with $C_r$ independent
of $q$. For two active incidences, the $O_r(q^{-1})$ conclusion holds
already for $r\ge2$. Thus, in the notation of the source, the estimate
labelled \texttt{eq:transverse-local-l1} can be replaced by
$O_r(q^{-1})$ when $r\ge3$.
\end{theorem}
\begin{proof}
We first establish the algebraic parametrization. If $B$ is nonsingular
in $\Sym_{2r}$ and $A_0$ has rank $r$, then
\begin{equation}\label{eq:projection-identity}
 \rank(B-A_0)=r\quad\Longleftrightarrow\quad A_0B^{-1}A_0=A_0.
\end{equation}
If the two ranks are $r$, their images are complementary because their
sum is $B$. The decomposition
$A_0B^{-1}+(B-A_0)B^{-1}=I$ therefore consists of the two complementary
projections, which proves the forward direction. Conversely,
$A_0B^{-1}$ is an idempotent of rank $r$ under the displayed identity,
and its complementary idempotent has rank $r$, proving the reverse.
If $A_0$ and $A_0-B$ have quotient signs $\xi_1,\xi_2$, their
complementary image representation gives
\begin{equation}\label{eq:pair-match}
 \chi(\det B)=\eps^r\xi_1\xi_2.
\end{equation}
Indeed $B$ is congruent to $R_1\oplus(-R_2)$ in bases of those images.

Suppose first that all three incidences are active. On $\mathcal T$ put
\[
 B_0=X_1-X_2,\quad B_1=X_1-X_3,\quad
 Q=B_0^{-1},\quad L=B_0^{-1}-B_1^{-1},\quad A_0=X_1-Y.
\]
The identity $L=B_0^{-1}(B_1-B_0)B_1^{-1}$ shows that $L$ is symmetric
and nonsingular. Write $A_0=PRP^{\mathsf T}$, where the columns of $P$
span its $r$-dimensional image. Equation~\eqref{eq:projection-identity}
for $B_0$ and $B_1$ is equivalent to
\begin{equation}\label{eq:local-parametrization}
 P^{\mathsf T}LP=0,\qquad P^{\mathsf T}QP\text{ nonsingular},\qquad
 R=(P^{\mathsf T}QP)^{-1}.
\end{equation}
Consequently the centers satisfying all three rank conditions are in
bijection with maximal $L$-isotropic subspaces on which $Q$ restricts
nonsingularly. The image recovers the subspace and the last equation
recovers the center, so there is no multiplicity.

Let $\mathcal M$ be the event that every pair match
\eqref{eq:pair-match} holds. If a match fails, both $F_A$ and $G_A$
vanish. If all matches hold, multiplication of their three determinant
characters gives $\chi(\det L)=\eps^r$, and hence $L$ is split.
Here is an elementary justification of this last criterion. A diagonal
nondegenerate form in dimension at least three has a nonzero isotropic
vector: for nonzero $a,b,c$, the two sets $\{ax^2:x\in\F_q\}$ and
$\{-c-by^2:y\in\F_q\}$ have size $(q+1)/2$ and intersect. Split off
hyperbolic planes until the anisotropic remainder has dimension at most
two. An anisotropic binary form has determinant character $-\eps$,
whereas a hyperbolic plane has character $\eps$. Thus a nondegenerate
form of dimension $2r$ is split exactly when its determinant character
is $\eps^r$.

Inversion in \eqref{eq:local-parametrization} preserves nonzero
determinant character. Once the sign on the first image equals
$\xi_1$, the pair matches force the other two signs as well. If
$f_{\xi_1}(Q;L)$ denotes the restriction-sign fraction among the
$N_r$ generators of split $L$, the exact normalized count on $\mathcal M$
is therefore
\begin{equation}\label{eq:local-triple-factor}
 F_A=\frac{q^{-c(D)}N_r}
 {\lambda_{r,\xi_1}\lambda_{r,\xi_2}\lambda_{r,\xi_3}}
 f_{\xi_1}(Q;L)
 =(16+O_r(q^{-1}))f_{\xi_1}(Q;L).
\end{equation}
The leading coefficient follows from
$\lambda_{r,\xi}=(1/2+O_r(q^{-1}))q^{-c(r)}$ and
$r(r-1)/2-c(2r)+3c(r)=0$.

It remains to justify the probability law to which the averaging theorem
is applied. Before invertibility conditioning, $B_0,B_1$ are independent
uniform symmetric matrices, by conditioning on $X_1$ and translating
$X_2,X_3$. Conditional on both being nonsingular, inversion preserves
the independent uniform laws on nonsingular symmetric matrices.
Conditional further on $L=Q-(Q-L)$, $Q$ is uniform on
$\{Q:Q,Q-L\text{ nonsingular}\}$: each such $Q$ gives exactly one
ordered inverse pair. The third transverse condition is precisely
nonsingularity of $L$. For every fixed nonsingular split $L$,
Theorem~\ref{thm:polar} and Lemma~\ref{lem:conditioning} therefore give
\[
 \E\bigl[\one_{\mathcal M}|f_{\xi_1}(Q;L)-1/2|
       \mid\mathcal T,L\bigr]=O(q^{-1})\qquad(r\ge3).
\]
The matching indicator stays \emph{inside} this expectation; there is
no unjustified conditioning on the matches. For nonsplit $L$, the match
event is empty. Combining this bound with
\eqref{eq:local-triple-factor}, whose target is $8$ on a match, and
integrating over $L$ proves \eqref{eq:local-improvement} for three active
incidences. Since $0\le f_{\xi_1}\le1$, it also proves the uniform bound
on $F_A$ there.

For two active incidences, relabel them $1,2$ and set $B=X_1-X_2$,
$Q=B^{-1}$. The same projection identity parametrizes centers by all
$r$-subspaces on which $Q$ is nonsingular. On the pair match the exact
count is
\begin{equation}\label{eq:local-pair-factor}
 F_A=\frac{q^{-c(D)}\gb Dr}
 {\lambda_{r,\xi_1}\lambda_{r,\xi_2}}
 \overline\I_{\xi_1}^{\,\mathrm{Gr}}(Q)
 =(4+O_r(q^{-1}))\overline\I_{\xi_1}^{\,\mathrm{Gr}}(Q),
\end{equation}
and off the match both quantities vanish. Conditional on $B$ being
nonsingular, $Q$ is uniform nonsingular symmetric. Theorem~\ref{thm:grassmann}
and Lemma~\ref{lem:conditioning}, with the match indicator retained
inside the expectation, give error $O_r(q^{-1})$ for $r\ge2$.
Restricting the nonnegative error to the smaller event $\mathcal T$
can only decrease its expectation. Formula~\eqref{eq:local-pair-factor}
also gives the uniform bound. For zero or one active incidence,
$F_A=G_A=1$ by normalization and translation invariance.
\end{proof}

\begin{remark}[Scope of the application]
The statement just proved is a local replacement in an explicitly
identified source argument. It does not analyze the other sections of
the Sidorenko manuscript, improve its global quantitative choices, or
certify its claimed counterexample. In particular, a better local error
need not improve a global theorem if another error term dominates.
The two-family exact identities are unconditional identities; their
conditional upper bounds follow from Lemma~\ref{lem:conditioning},
not from a claim that conditioning preserves the entire Fourier law.
\end{remark}

\section{Exact checks, examples, and reproducibility}
\label{sec:checks}
The accompanying programs separate closed-form rational evaluation from direct
finite-field enumeration. The module \texttt{code/exact\_formulas.py} implements
the marginal laws, extension probabilities, covariance kernels, intersection
profiles, and common-generator moments using Python integers and
\texttt{fractions.Fraction}. The program \texttt{code/verify.py} then compares
these formulas with explicit matrices and subspaces. Every pass/fail comparison
is exact. Numerical evidence supports the audit; the proofs in the preceding
sections, not the finite tests, establish the general theorems.

\subsection{An exhaustive four-dimensional example}
Take $q=3$, $r=2$, and the standard hyperbolic form on $\F_3^4$.
There are $3^{10}=59\,049$ symmetric forms $Q$ and eight generators, four in
each family. The verifier enumerates all of them. It obtains
\begin{align*}
 \E\overline\J&=\frac23,&
 \E\overline\D&=-\frac29,&
 \E\overline\I_+&=\frac29,&
 \E\overline\I_-&=\frac49,\\
 \Var(\overline\J)&=\frac1{36},&
 \Var(\overline\D)&=\frac{41}{324},&
 \Cov(\overline\J,\overline\D)&=-\frac1{108}.
\end{align*}
The contrast covariance matrix is
\[
 \begin{pmatrix}
 \E\Delta_\J^2&\E\Delta_\J\Delta_\D\\
 \E\Delta_\J\Delta_\D&\E\Delta_\D^2
 \end{pmatrix}
 =\begin{pmatrix}1/9&-1/27\\-1/27&1/9\end{pmatrix}.
\]
Consequently
$\E\Delta_{\I_+}^{\,2}=1/27$ and
$\E\Delta_{\I_-}^{\,2}=2/27$, in agreement with
Theorem~\ref{thm:contrast}. Notice that neither individual-sign mean equals
$1/2$ in this example. Centering at the exact mean is essential.

For the signed common-generator count $Z$, the complete histogram is
\begin{center}
\begin{tabular}{@{}lrrrrr@{}}
\toprule
Value of $Z$ & $-2$ & $-1$ & $0$ & $1$ & $2$\\
\midrule
Number of forms &108&5\,400&48\,033&5\,400&108\\
\bottomrule
\end{tabular}
\end{center}
These counts agree with Theorem~\ref{thm:rank-two-law}:
$a_3=200/2187$, $b_3=4/2187$, and
$\Pp(Z=0)=593/729$. The pointwise identity is
$\Delta_\D=3Z/4$, since $\chi(-1)=-1$ in $\F_3$.
The characteristic polynomial obstruction was checked for every one of the
$13\,635$ forms having at least one common generator, using polynomial
coefficient arithmetic rather than Newton identities that would divide by
zero in characteristic three.

\subsection{Coverage and limitations of the checks}
The closed-form suite checks 96 pairs $(r,q)$ with $1\le r\le16$ and
$q\in\{3,5,7,9,11,25\}$, as well as 120 Grassmannian parameter triples.
The appearances of $q=9,25$ here evaluate exact prime-power counting formulas;
\emph{they do not perform field arithmetic modulo 9 or 25}.
All direct field enumerations use $\F_3$.

The direct suite enumerates every symmetric matrix in dimensions one, two,
and three and checks every nonsingular Fourier frequency in those dimensions.
The Fourier check works in $\mathbb Z[\zeta_3]$: if the three residue-class
counts are $c_0,c_1,c_2$, its squared norm is the integer
$c_0^2+c_1^2+c_2^2-c_0c_1-c_0c_2-c_1c_2$.
It also checks all conditional extensions of every leading block of sizes
one and two, and all four joint-sign probabilities for
$(r,s,t)=(2,2,0),(2,2,1),(2,3,1),(2,2,2)$.

A separate odd-rank pointwise check uses all 80 generators of a split form
on $\F_3^6$ and every one of the $3^9=19\,683$ matrices in the slice
\[
 Q=\begin{pmatrix}0&B\\B^{\mathsf T}&C\end{pmatrix},
 \qquad B\text{ diagonal},\quad C\in\Sym_3(\F_3).
\]
This verifies the identity for $\Delta_\J$ on that slice. It is not an
enumeration of all $3^{21}$ forms in $\Sym_6(\F_3)$.

The actual successful transcript is shipped as
\texttt{verification-results.txt}. The programs require Python with NumPy;
run \texttt{python code/verify.py} without the \texttt{-O} option, so that
assertions remain enabled. No randomized testing, floating-point tolerance,
proof-assistant certification, or external referee review is represented by
this transcript.

\section{Further research questions and topics}
\label{sec:questions}
The following are specific research directions suggested by the proved
results. They are proposed here; this list is not a claim that every question
has appeared previously as a named open problem.

\begin{question}[Self-adjoint pencils and the distribution of common generators]
Classify the possible values of $Z(Q)$ in terms of the primary decomposition
of the self-adjoint operator $H^{-1}Q$, including nonsemisimple blocks.
Determine the full joint distribution of $(C_+,C_-)$ for $r=3$ and beyond.
The square-characteristic-polynomial obstruction is necessary for a common
generator but is not, by itself, a classification of nonzero $Z$: the scalar
pencil $Q=aH$ has $Z=0$. A useful next result would distinguish zero counts,
balanced positive counts, and genuine family imbalance by module invariants.
\end{question}

\begin{question}[The remaining parity and higher moments]
The variance-normalized contrast has a degenerate limit for even-rank
$\D$ and odd-rank $\J$, by Corollary~\ref{cor:non-gaussian}.
What are the laws for odd-rank $\D$ and even-rank $\J$? Determine higher
moments in all parities, their dependence on $q$ and $r$, and the law
conditional on a nonzero contrast. The rare-event theorem already excludes
a nondegenerate Gaussian limit in the two parities it treats; a conjecture
about the other parities should not ignore that obstruction.
\end{question}

\begin{question}[Ordinary averaging and simultaneous asymptotic regimes]
Find sharp higher moments and a distributional description for
$\overline\D-\mu_r$ and $\overline\I_\xi-\pi_{r,\xi}$.
Do different behaviors occur for fixed $r$ with $q\to\infty$, fixed $q$
with $r\to\infty$, and jointly growing parameters?
The exact covariance kernels and incidence profiles give second moments in
all these regimes, but fixed-rank $O_r(\cdot)$ estimates alone do not give a
uniform distributional theorem. Identify which geometric degeneracies
carry the ordinary variance, rather than assuming independent sampling.
\end{question}

\begin{question}[Three or more restriction laws]
For two subspaces, their dimensions and intersection dimension suffice for
the joint law. For three or more subspaces, identify the additional
configuration invariants needed for an exact analogue of
Theorem~\ref{thm:joint}. Derive a practical recursive calculus that keeps
singular overlaps and compatibility constraints without treating the
restrictions as conditionally independent when they are not. This is a
natural route to higher moments and to face systems sharing several point
matrices in rank-layer constructions.
\end{question}

\begin{question}[Smaller signed designs and partial spectral cancellation]
Theorem~\ref{thm:filter-rigidity} classifies universal cancellation of every
rank below $r$ when all generators are available: only the Latin--Greek
signed weight survives, up to scaling. What are optimal support sizes and
weight norms when cancellation is required only through rank $s<r$?
Can a small signed subfamily reproduce a prescribed part of the discriminant
contrast, or can one prove a support-size obstruction from the incidence
matrix? The exact energy identity provides an objective function for this
optimization rather than merely a heuristic for sampling.
\end{question}

\begin{question}[Other polar geometries and characteristic two]
Develop analogous incidence-to-Fourier identities for Hermitian or
alternating forms, and identify the statistic replacing determinant
character. In characteristic two, symmetric bilinear and quadratic forms
have different structural features; the trace pairing and type invariants
used here cannot simply be copied. The association-scheme framework of
\cite{schmidt2020} is a starting point, but a useful extension should state
its Fourier dual space and geometric design hypothesis explicitly.
\end{question}

\begin{question}[Sharp finite-field bounds and conditioning]
Optimize the uniform constants in the covariance and averaging inequalities,
and compute exact conditional moments for a fixed finite pencil of
invertibility constraints. Lemma~\ref{lem:pencil-conditioning} gives a
relative $1+O_{r,k}(q^{-1})$ comparison because the observables are constant
along pencil lines. Can the first correction term be expressed using the
characteristic polynomial and the restriction profile? This would separate
an exact conditioning effect from the deterministic bias $\pi_{r,\xi}-1/2$.
\end{question}

\begin{question}[Propagation through the rank-layer construction]
Track Theorem~\ref{thm:local} through every other error estimate in the
cited transverse and rank-layer arguments. Determine whether the improved
local rate changes an admissible field-size or dimension choice once the
other sources of error are retained. A full answer must audit those other
inputs; an $O(q^{-1})$ replacement in one local lemma is not automatically
an improved global counterexample parameter. It may also be useful to
retain the two families separately and investigate whether their exact
agreement off the rare pencil locus simplifies a signed construction.
\end{question}

\begin{question}[A modular formalization]
Formalize, in order: the extension law and finite counting identities;
the two-subspace compatibility map; the weighted Fourier-transfer theorem;
the two-family incidence cancellation; and the pointwise rare-event identity.
The first two stages use finite linear algebra and finite sums only. The
geometric parity argument and the additive-character normalization should
be separate reusable modules. Finally connect the local rank-layer theorem
to the source's precise definitions rather than assigning formal status
from its inclusion in a research directory. The supplied code is a test
suite, not a substitute for any of these proof obligations.
\end{question}

\appendix
\section{Pinned sources and integration scope}
The repository versions used for the final source crosswalk are:
\begin{center}
\begin{tabular}{@{}ll@{}}
\toprule
Repository & Commit\\
\midrule
\texttt{openai/math} & \texttt{adc7f1241b42e322a6451854ab7e4b4c146bf78a}\\
\texttt{VladimirReshetnikov/ProveIt} & \texttt{a41aa448d67395387593f1931cda142836afccdc}\\
\bottomrule
\end{tabular}
\end{center}
The direct OpenAI source is the file \texttt{transverse.tex} in the
September-23-2026 Sidorenko manuscript's \texttt{build/sections/} directory.
Its exact blob is
\begin{center}
\texttt{d4991976e630f117c8d85a7e22deb60643a2450d}.
\end{center}
The reviewed source labels are
\begin{center}
\begin{tabular}{l}
\texttt{lem:transverse-restrictions}\\
\texttt{lem:local-transverse}\\
\texttt{eq:transverse-three-normalization}\\
\texttt{eq:transverse-two-normalization}.
\end{tabular}
\end{center}
In that source, $b(n)$ denotes what this article calls $c(n)$.
The full repository path is recorded in the accompanying source audit.

The ProveIt inventory is the \texttt{README.md} in the research directory
whose path components are
\begin{center}
\begin{tabular}{l}
\texttt{Combinatorics/Ramsey/Research/GowersSzemeredi/}\\
\texttt{local-quantitative-refinements/}.
\end{tabular}
\end{center}
Its blob is \texttt{02f1eb26decc04804b38d00bee586568701ff88e}.
The selected companion audit and its blob are
\begin{center}
\begin{tabular}{l}
\texttt{33-arrangement-stability-SOURCE\_AUDIT.md}\\
\texttt{9ed4c250339f2f6172ff0765adcc3e34e24b699b}.
\end{tabular}
\end{center}
They establish context and claim boundaries, not premises for the theorems
proved here. More detailed provenance and a theorem-status ledger accompany
the source archive.

No repository was modified during preparation. The deliverable contains
this manuscript, its PDF, the verification programs and their transcript,
and integration notes. It does not redistribute the cited third-party
manuscripts or imply their authors' endorsement. Proposed inclusion in
ProveIt does not by itself confer formal or independently reviewed status.

\clearpage
\begin{thebibliography}{9}
\bibitem{openai}
\texttt{openai/math} contributors,
\emph{A counterexample to Sidorenko's conjecture}, repository manuscript
filed under 23 September 2026, section ``Transverse configurations and
determinant signs.''
\href{https://github.com/openai/math/tree/adc7f1241b42e322a6451854ab7e4b4c146bf78a/preprints/A-counterexample-to-Sidorenkos-conjecture-September-23-2026}{Pinned repository source}.
Accessed 7 October 2026; exact file and blob are recorded in Appendix A.
This citation attributes the source artifact, not independent verification
of its global claim.

\bibitem{proveit}
\texttt{VladimirReshetnikov/ProveIt},
\emph{Local Quantitative Refinements of Gowers's Proof of Szemer\'edi's
Theorem}, research inventory and
\texttt{33-arrangement-stability-SOURCE\_AUDIT.md}.
\href{https://github.com/VladimirReshetnikov/ProveIt/tree/a41aa448d67395387593f1931cda142836afccdc/Combinatorics/Ramsey/Research/GowersSzemeredi/local-quantitative-refinements}{Pinned research directory}.
Accessed 7 October 2026. Used for research context and provenance, not as
an unverified theorem input.

\bibitem{schmidt2015}
K.-U. Schmidt,
\emph{Symmetric bilinear forms over finite fields with applications to
coding theory}, Journal of Algebraic Combinatorics \textbf{42} (2015),
635--670. \href{https://arxiv.org/abs/1410.7184}{arXiv:1410.7184}.

\bibitem{schmidt2020}
K.-U. Schmidt,
\emph{Quadratic and symmetric bilinear forms over finite fields and their
association schemes}, Algebraic Combinatorics \textbf{3} (2020), 161--189.
\href{https://doi.org/10.5802/alco.88}{doi:10.5802/alco.88}.

\bibitem{polar-designs}
M. Kiermaier, K.-U. Schmidt and A. Wassermann,
\emph{Designs in finite classical polar spaces}, Designs, Codes and
Cryptography \textbf{93} (2025), 1143--1162.
\href{https://doi.org/10.1007/s10623-024-01491-x}{doi:10.1007/s10623-024-01491-x}.
Published online 17 September 2024. See the Latin--Greek halving discussion
and Lemma 14.
\end{thebibliography}
\end{document}
